% !TeX root = ../../Skript.tex
\cohead{\Large\textbf{Funktionsbegriff}}
\fakesubsection{Funktionsbegriff}
\begin{tcolorbox}
    Funktion:

    \bigskip

    \textcolor{loestc}{Eine Funktion ist eine Zuordnung, die \textbf{jedem} x-Wert aus der Definitionsmenge \textbf{genau einen} y-Wert zuordnet.
    }

    \bigskip

    \textcolor{loestc}{\(f:\ D_f\rightarrow W_f\)}
\end{tcolorbox}
\(f\): \textcolor{loes}{Name der Funktion. Funktionen werden normalerweise mit Kleinbuchstaben bezeichnet.}

\medskip

\(D_f\): \textcolor{loes}{Definitionsmenge der Funktion \(f\). Die Definitionsmenge umfasst alle Werte, die man für \(x\) einsetzen darf.}

\medskip

\(W_f\): \textcolor{loes}{Wertemenge der Funktion \(f\). Alle möglichen Werte, die die Funktion annehmen kann, d.h. für jedes Element \(y\) der Wertemenge gibt es auch ein Element \(x\) aus der Definitionsmenge, so dass die Funktion den Wert \(y\) annimmt, wenn man \(x\) einsetzt.}

Beispiel: \(y=x^2\)

\textcolor{loes}{Diese Funktion ordnet jedem x-Wert das Quadrat dieses Wertes als y-Wert zu, d.h. der 1 wird die 1 zugeordnet, der 2 die 4, usw.}

\textcolor{loes}{Wählt man \(D_f=\R\), so ergibt sich \(W_f=\R^0_+=\{y\in\R\vert y\geq 0\}\)}

\bigskip

\begin{tcolorbox}
    Schreibweise für Funktionsgleichungen:

    \bigskip

    \textcolor{loestc}{Bisher wurden Funktionen z.B. als \(y=2x+1\) geschrieben.}

    \bigskip

    \textcolor{loestc}{Neue Schreibweise der Funktionsgleichung: \(f(x)=2x+1\)}

    \bigskip

    \textcolor{loestc}{\(f\): Name der Funktion, in den meisten Fällen \(f,\ g,\ h\) oder \(f_1,\ ,f_2,\ f_3\) usw.}

    \bigskip

    \textcolor{loestc}{\(x\): Laufvariable, wird bei uns so gut wie immer mit \(x\) bezeichnet.}

    \bigskip

    \textcolor{loestc}{\(2x+1\): Funktionsterm}

    \bigskip

    \textcolor{loestc}{Gesprochen als \(f\) von \(x\) gleich Zwei \(x\) Plus Eins.}
\end{tcolorbox}
Der Vorteil der neuen Schreibweise ist, dass man Funktionen an Hand des Namens unterscheiden kann. Zudem erlaubt sie uns für \(x\) Werte einzusetzen und die zugeordneten y-Werte allgemeiner aufzuschreiben:
\begin{tcolorbox}
    Einsetzen von Werten:

    \bigskip

    \textcolor{loestc}{Für \(x\) können nun verschiedene Werte eingesetzt werden:, z.B.:}

    \bigskip

    \textcolor{loestc}{\(x=5\): \(f(5)=2\cdot 5+1=11\)}

    \bigskip

    \textcolor{loestc}{Gesprochen: \(f\) von 5 ist gleich 11}

    \bigskip

    \textcolor{loestc}{\(f(5)\) bezeichnet man auch als Funktionswert von \(f\) an der Stelle 5.}

    \textcolor{loestc}{Grafisch bedeutet dies, dass das Schaubild von \(f(x)\) durch den Punkt \(P(5|11)\) verläuft.}
\end{tcolorbox}
\newpage
%%%%%%%%%%%%%%%%%%%%%%%%%%%%%%%%%%%%%%%%%%%%%%%%%%%%%%%%%%%%%%%%%%%%%%%%%%%%%%%%%%%%%%%%%%%%%%%%%%%%%%%
\begin{Exercise}[title={\raggedright Finde die passenden Paare gleichwertiger Aussagen.}, label=fktBegriffA1]

    \begin{minipage}{\textwidth}
        \adjustbox{valign=t}{\begin{minipage}{0.5\linewidth}
                \begin{enumerate}[label=\alph*)]
                    \item Der Funktionswert an der Stelle 3 ist 4.
                    \item \(f\) von 8 ist 0.
                    \item \(g(x)=0,5x\)
                    \item \(P(-5|1)\) liegt auf dem Schaubild der Funktion.
                    \item Die Funktion ordnet jedem \(x\) das Dreifache des Wertes von \(x\) zu.
                    \item Der Punkt \(P(-1|4)\) liegt auf dem Schaubild der Funktion.
                    \item Die Funktion ordnet der 0 die 5 zu.
                    \item \(f(4)=3\)
                    \item Der Punkt \(Q(0|0)\) liegt auf dem Schaubild der Funktion.
                \end{enumerate}
        \end{minipage}}%
        \adjustbox{valign=t}{\begin{minipage}{0.5\linewidth}
                \begin{enumerate}[label=\arabic*)]
                    \item \(f(-5)=1\)
                    \item Das Schaubild der Funktion verläuft durch den Ursprung.
                    \item An der Stelle 4 ist der Funktionswert 3
                    \item \(h(-1)=4\)
                    \item \(f_2(0)=5\)
                    \item \(f(8)=0\)
                    \item Die Funktionsgleichung lautet \(g(x)=3x\)
                    \item \(f(3)=4\)
                    \item Die Funktion ordnet jedem  \(x\) den halben Wert als Funktionswert zu.
                \end{enumerate}
        \end{minipage}}
    \end{minipage}
\end{Exercise}
\begin{Exercise}[title={\raggedright Entscheide jeweils, ob das Schaubild zu einer Funktion gehört und begründe deine Entscheidung.}, label=fktBegriffA2]

\begin{minipage}{\textwidth}
    \adjustbox{valign=t, padding=0ex 0ex 1ex 0ex}{\begin{minipage}{0.25\linewidth}
        \includegraphics[width=\textwidth]{\fktBegriff/pics/GeschwindigkeitMassepunkt.png}
    \end{minipage}}%
    \adjustbox{valign=t, padding=1ex 0ex 1ex 0ex}{\begin{minipage}{0.25\linewidth}
        \includegraphics[width=\textwidth]{\fktBegriff/pics/xsinx.png}
    \end{minipage}}%
    \adjustbox{valign=t, padding=1ex 0ex 1ex 0ex}{\begin{minipage}{0.25\linewidth}
            \includegraphics[width=\textwidth]{\fktBegriff/pics/xcosxGedreht.png}
    \end{minipage}}%
    \adjustbox{valign=t, padding=1ex 0ex 0ex 0ex}{\begin{minipage}{0.25\linewidth}
            \includegraphics[width=\textwidth]{\fktBegriff/pics/Kreis.png}
    \end{minipage}}%
\end{minipage}
\end{Exercise}
%%%%%%%%%%%%%%%%%%%%%%%%%%%%%%%%%%%%%%%%%
\begin{Answer}[ref=fktBegriffA1]

    \begin{enumerate}[label=\alph*)]
        \item 8)
        \item 6)
        \item 9)
        \item 1)
        \item 7)
        \item 4)
        \item 5)
        \item 3)
        \item 2)
    \end{enumerate}
\end{Answer}
%%%%%%%%%%%%%%%%%%%%%%%%%%%%%%%%%%%%%%%%%%%%%%%%%%%%%%%%%%%%%%%%%%%%%%%%%%%%%%%%%%%%%%%%%%%%%%%%%%%%%%%
%%%%%%%%%%%%%%%%%%%%%%%%%%%%%%%%%%%%%%%%%
\begin{Answer}[ref=fktBegriffA2]

    \begin{enumerate}[label=\alph*)]
        \item Schaubild: Keine Funktion, da dem \(x\)-Wert 5 unendlich viele \(y\)-Werte zugeordnet werden.
        \item Schaubild: Funktion, da jedem
    \end{enumerate}
\end{Answer}